% ============================================================
% MODA*-G: A Robust Mixed-Data Outlier Detection Framework
%          with Attention-based Softmax Gating
%
% ArXiv submission — stat.ME / stat.CO
% Author: Luis J. Marcano-Verde
% Version: FINAL — March 2026
%
% Changes from draft:
%  - Author info updated (email, affiliation)
%  - Condition (C3) replaced by (C3') — mixture model
%  - Condition (C5) added — finite fourth moments
%  - Proposition 2 added — O_P(1/sqrt(n)) convergence rate
%  - Section 8 (Limitations) added
%  - Related Work expanded with Marcano & Fermin (2013)
%  - Abstract updated: "five theorems and two propositions"
%  - Calibration constant c specified in SDG_mad definition
%  - Corollary numbering fixed
% ============================================================

\documentclass[11pt]{article}

% ── Packages ──────────────────────────────────────────────
\usepackage[margin=1in]{geometry}
\usepackage{amsmath,amssymb,amsthm}
\usepackage{mathtools}
\usepackage{bm}
\usepackage{booktabs}
\usepackage{array}
\usepackage{multirow}
\usepackage{xcolor}
\usepackage{hyperref}
\usepackage{natbib}
\usepackage{graphicx}
\usepackage{caption}
\usepackage{subcaption}
\usepackage{algorithm}
\usepackage{algpseudocode}
\usepackage{microtype}
\usepackage{setspace}
\usepackage{enumitem}
\usepackage{thmtools}

% ── Colors ────────────────────────────────────────────────
\definecolor{navy}{RGB}{26,58,92}
\definecolor{teal}{RGB}{14,102,85}
\definecolor{crimson}{RGB}{192,57,43}
\definecolor{gold}{RGB}{184,134,11}

% ── Theorem environments ──────────────────────────────────
\theoremstyle{plain}
\newtheorem{theorem}{Theorem}
\newtheorem{proposition}{Proposition}
\newtheorem{corollary}{Corollary}
\newtheorem{lemma}{Lemma}

\theoremstyle{definition}
\newtheorem{definition}{Definition}
\newtheorem{condition}{Condition}
\newtheorem{remark}{Remark}

% ── Math operators ────────────────────────────────────────
\DeclareMathOperator{\MAD}{MAD}
\DeclareMathOperator{\median}{median}
\DeclareMathOperator{\tr}{tr}
\DeclareMathOperator{\diag}{diag}
\DeclareMathOperator*{\argmin}{arg\,min}
\DeclareMathOperator*{\argmax}{arg\,max}
\newcommand{\R}{\mathbb{R}}
\newcommand{\C}{\mathcal{C}}
\newcommand{\X}{\mathcal{X}}
\newcommand{\E}{\mathbb{E}}
\newcommand{\Prob}{\mathbb{P}}
\newcommand{\bx}{\mathbf{x}}
\newcommand{\bS}{\mathbf{S}}
\newcommand{\bw}{\mathbf{w}}
\newcommand{\bp}{\mathbf{p}}
\newcommand{\bone}{\mathbf{1}}
\newcommand{\inprob}{\xrightarrow{P}}
\newcommand{\indist}{\xrightarrow{d}}
\newcommand{\moda}{\text{MODA}^*\text{-G}}

% ── Hyperref ──────────────────────────────────────────────
\hypersetup{
  colorlinks=true,
  linkcolor=navy,
  citecolor=teal,
  urlcolor=teal,
  pdftitle={MODA*-G: A Robust Mixed-Data Outlier Detection
    Framework with Attention-based Softmax Gating},
  pdfauthor={Luis J. Marcano-Verde},
  pdfsubject={DOI: 10.5281/zenodo.19433757}
}

% ── Title ─────────────────────────────────────────────────
\title{
  \textbf{MODA*-G: A Robust Mixed-Data Outlier Detection
  Framework with\\Attention-based Softmax Gating}
}

\author{
  Luis J.\ Marcano-Verde\thanks{
    Correspondence: \href{mailto:luisjmarcanoverde@gmail.com}
    {luisjmarcanoverde@gmail.com}}\\[4pt]
  \small Independent Statistical Consultant and Data Scientist\\
  \small New York, NY, USA
}

\date{March 2026\\[4pt]
  \small\texttt{DOI: \href{https://doi.org/10.5281/zenodo.19433757}{10.5281/zenodo.19433757}}}

% ============================================================
\begin{document}
\maketitle

% ── Abstract ──────────────────────────────────────────────
\begin{abstract}
We introduce \textbf{MODA*-G} (\emph{Mixed-Data Outlier
Divergence Analysis with Softmax Gating}), a novel
unsupervised anomaly detection framework for tabular data
containing both numerical and categorical variables.
MODA*-G combines four complementary statistical
engines---a robust MCD-based Mahalanobis distance
($\text{SDN}_\text{mcd}$), a directional kurtosis score
with robust MAD evaluation ($\text{SDN}_\text{pe\~{n}a}$),
an original Gower-MAD distance score
($\text{SDG}_\text{mad}$), and a categorical entropy score
($\text{SDC}$)---through a Softmax Gating mechanism with
temperature parameter $k$, yielding observation-level
adaptive weights without heuristic rules.

We establish five formal theorems (consistency, 50\%
breakdown point, partial affine equivariance, score
comparability, and stochastic dominance) and two
convergence propositions: one establishing that the gating
weight of the dominant engine converges to one as
$k \to \infty$, and one establishing that the MODA*-G
score converges at rate $O_P(n^{-1/2})$.

Monte Carlo simulation ($n{=}300$, $p{=}4$, $q{=}2$,
100 replications per cell) across four canonical
contamination types confirms that MODA*-G with $k{=}3$
achieves mean AUC-ROC of 0.978, outperforming Isolation
Forest (0.911) and LOF (0.623) in all four scenarios
simultaneously. On the Annthyroid real-data benchmark
\citep{campos2016} ($n{=}7{,}200$, $p{=}6$, $q{=}15$),
MODA*-G achieves AUC-ROC $= 0.939$, outperforming IF
($0.632$), EIF ($0.496$), COPOD ($0.553$), CatBoost-AD
($0.554$), LOF ($0.620$), and surpassing the best
result of \citet{campos2016} by $+0.189$ AUC.
The framework operates natively in
$\X = \R^p \times \C^q$ without encoding categorical
variables, and produces an interpretable per-engine
diagnostic decomposition. Applications to epidemiological
biosurveillance are discussed.
\end{abstract}

\noindent\textbf{Keywords:} mixed data, outlier detection,
robust statistics, minimum covariance determinant, kurtosis
projection pursuit, Gower distance, softmax attention,
breakdown point, biosurveillance.

\noindent\textbf{MSC 2020:} 62G35, 62H15, 62P10.

% ============================================================
\section{Introduction}
\label{sec:intro}
% ============================================================

Anomaly detection in mixed-type tabular data---combining
continuous numerical measurements with unordered categorical
variables---remains an open methodological challenge.
Real-world data systems are almost universally mixed:
electronic health records integrate laboratory values with
ICD-10 diagnosis codes; epidemiological surveillance
databases combine incidence rates with pathogen identifiers
and geographic classifications; official statistical
microdata couple continuous indicators with categorical
sector codes. In each setting, outlier detection requires
methods that respect the distinct statistical natures of
numerical and categorical variables.

Existing approaches address this problem inadequately.
Classical multivariate outlier detectors---the Minimum
Covariance Determinant (MCD) estimator
\citep{rousseeuw1985}, the robust Mahalanobis distance
\citep{rousseeuw1990}, and the directional kurtosis of
\citet{pena2001}---operate exclusively in $\R^p$ and
require artificial encoding of categorical variables,
distorting the metric structure. Machine learning methods
including Isolation Forest \citep{liu2008} and Local
Outlier Factor \citep{breunig2000} inherit this limitation
and provide no formal statistical guarantees. Recent
mixed-data proposals \citep{yuan2023, cabral2025, wang2021}
use internal conversion mechanisms and lack formally
established properties.

The empirical foundation for the present work is
\citet{marcano2013}, who demonstrated through a simulation
study comparing seven multivariate outlier detection
methods that no single procedure dominates across all
contamination scenarios: the Curtosis-1 procedure
\citep{pena2001} detected as few as 3.58\% of true outliers
under correlated contamination at 20\% contamination rate,
while performing well in other settings. This instability
across scenarios directly motivates a multi-engine framework
combining complementary detection strategies with formal
statistical guarantees.

\paragraph{Contributions.}
\begin{enumerate}[label=(\roman*)]
  \item \textbf{Gower-MAD}: an original robustification
    of the Gower distance coefficient \citep{gower1971}
    replacing the classical range denominator with
    $1.4826 \cdot \MAD$, raising the breakdown point from
    0\% to 50\%.
  \item \textbf{Robust Pe\~{n}a evaluation}: an original
    modification of the Curtosis-1 projection
    \citep{pena2001} evaluating projections via median/MAD
    instead of classical moments, raising the breakdown
    point from 0\% to 50\%.
  \item \textbf{MODA*-G}: a unified mixed-data anomaly
    detection framework combining four complementary
    engines through Softmax Gating---an attention mechanism
    that adaptively amplifies the dominant engine at the
    observation level without heuristic rules.
  \item \textbf{Formal guarantees}: five theorems and two
    propositions, providing the first formal statistical
    properties---including 50\% breakdown point and
    $O_P(n^{-1/2})$ convergence rate---for a mixed-data
    outlier detection method.
  \item \textbf{Empirical superiority}: Monte Carlo evidence
    that MODA*-G with $k=3$ outperforms Isolation Forest
    in all four canonical contamination scenarios (mean
    AUC-ROC 0.978 vs.\ 0.911).
\end{enumerate}

% ============================================================
\section{Related Work}
\label{sec:related}
% ============================================================

\paragraph{Comparative simulation studies.}
\citet{marcano2013} compared seven multivariate outlier
detection methods---including the Mahalanobis distance,
Curtosis-1 \citep{pena2001}, Minimum Covariance Determinant
\citep{rousseeuw1985}, and projection pursuit variants---
through Monte Carlo simulation across multiple contamination
structures and correlation scenarios. The study demonstrated
that detection performance is highly scenario-dependent and
that no single method achieves uniformly superior detection,
establishing the empirical need for a multi-engine framework.
The present work can be understood as a methodological
response to this finding: rather than selecting among
existing methods, MODA*-G combines them through a gating
mechanism that adapts to the contamination structure present
in each dataset.

\paragraph{Numerical outlier detection.}
The robust Mahalanobis distance based on the MCD estimator
\citep{rousseeuw1985, rousseeuw1999} achieves 50\%
breakdown point and detects isolated numerical outliers.
\citet{pena2001} proposed kurtosis-based projection pursuit
(Curtosis-1) specifically to detect clustered outliers that
cause masking in MCD-based methods, a complementary
capability that MODA*-G preserves through Engine II.
Isolation Forest \citep{liu2008} and LOF \citep{breunig2000}
are widely used but require encoding of categorical
variables and provide no formal breakdown point.

\paragraph{Mixed-data methods.}
WNIN \citep{wang2021} uses graph-based weights on mixed
variables but lacks formal breakdown point analysis and
requires internal encoding. PMIOD \citep{yuan2023} employs
mutual information weights but provides no formal
properties. HySortOD \citep{cabral2025} sorts mixed
variables internally. None of these methods produce
interpretable per-engine diagnostic decompositions.
MODA*-G differs by operating natively in $\R^p \times \C^q$,
providing formal guarantees, and enabling causal diagnosis
of each detected anomaly.

\paragraph{Attention mechanisms and mixture of experts.}
Softmax-weighted aggregation was introduced for neural
sequence modelling \citep{bahdanau2015} and formalized
in the Transformer architecture \citep{vaswani2017}.
Mixture-of-experts models \citep{jordan1994} use gating
functions to route inputs to specialized experts.
MODA*-G applies this paradigm to robust statistical
estimators for the first time, yielding self-normalizing
observation-level weights without external supervision
or labeled training data.

% ============================================================
\section{Data Space, Notation and Regularity Conditions}
\label{sec:setup}
% ============================================================

Let $\mathbf{X}_n = \{x_1, \ldots, x_n\}$ be an i.i.d.\
sample from distribution $P$ on the mixed product space
$\X = \R^p \times \C^q$,
where $\R^p$ is the $p$-dimensional Euclidean space and
$\C^q = \C_1 \times \cdots \times \C_q$ is the Cartesian
product of $q$ finite unordered categorical domains with
$|\C_j| = m_j$ categories. Each observation
$x = (x_\text{num}, x_\text{cat}) \in \X$ has a numerical
component $x_\text{num} \in \R^p$ and a categorical
component $x_\text{cat} \in \C^q$. No encoding of
categorical variables is performed at any stage.

We write $\median(\cdot)$ for the sample median,
$\MAD(\cdot) = \median(|X_i - \median(X)|)$ for the sample
median absolute deviation, and
$\widehat{\sigma}(\cdot) = 1.4826 \cdot \MAD(\cdot)$
for the robust scale estimator consistent for $\sigma$
under normality \citep{rousseeuw2018}.

\begin{definition}[Normalization function]
\label{def:phi}
For $d \geq 0$ and $\kappa > 0$:
$\varphi(d;\kappa) = d/(d + \kappa) \in [0,1)$.
With $\kappa = \sqrt{\chi^2_{p,0.975}}$, this maps the
robust Mahalanobis distance to $[0,1)$ such that 97.5\%
of observations from an elliptical null distribution
receive scores below 0.5.
\end{definition}

\noindent We assume the following regularity conditions.

\begin{condition}[Identifiability --- C1]
The marginal of $P$ on $\R^p$ is elliptically contoured
with unique location $T^* \in \R^p$ and positive definite
scatter $C^* \in \R^{p \times p}$.
\end{condition}

\begin{condition}[Clean majority --- C2]
The contamination fraction satisfies $\varepsilon <
1 - h/n$, where $h = \lfloor n/2 \rfloor + 1$.
\end{condition}

\begin{condition}[Generalised observability --- C3$'$]
\label{cond:c3prime}
The observed distribution follows a mixture model
$P = (1-\varepsilon)P_0 + \varepsilon P_\text{out}$,
where $P_0$ is the clean distribution with
$P_0(X_j = c_k) > 0$ for all observed categories
$c_k \in \mathcal{S}_j \subseteq \C_j$, and
$P_\text{out}$ is the contaminating distribution.
Categories in $\C_j \setminus \mathcal{S}_j$ may have
$P_0(X_j = c_k) = 0$ but $P_\text{out}(X_j = c_k) > 0$.
\end{condition}

\begin{remark}
\label{rem:c3}
Condition~(C3$'$) is more general than the standard
observability condition $P(X_j = c_k) > 0$ for all $c_k$.
It allows categories that appear exclusively in contaminated
observations---as in the O4 simulation scenario where
outliers carry category code 99, never observed in clean
data. Under (C3$'$) and the mixture model, the marginal
frequency of any category satisfies
$P(X_j = c) = (1-\varepsilon) P_0(X_j = c) +
\varepsilon P_\text{out}(X_j = c) > 0$
whenever $\varepsilon > 0$, restoring the applicability
of the SLLN to all category frequencies including
those of unseen-in-clean categories.
The Laplace-smoothed estimator
$\hat{P}_\text{unseen} = 1/(n + m_k + 1)$ converges
to $\varepsilon P_\text{out}(X_j = c)/n \to 0$ in
the limit, ensuring SDC correctly assigns high scores
to pure categorical outliers.
\end{remark}

\begin{condition}[Non-degeneracy --- C4]
All eigenvalues of $C^*$ are strictly positive.
\end{condition}

\begin{condition}[Finite fourth moments --- C5]
\label{cond:c5}
$\E[\|X_\text{num}\|^4] < \infty$ and all mixed moments
$\E[X_{\text{num},j}^2 X_{\text{num},k}^2] < \infty$
for $j, k = 1, \ldots, p$.
\end{condition}

\begin{remark}
Condition~(C5) ensures finite asymptotic variances for
all four engine scores, which is required for the
$O_P(n^{-1/2})$ convergence rate established in
Proposition~\ref{prop:rate}. Under normality of
$X_\text{num}$, (C5) is automatically satisfied.
\end{remark}

% ============================================================
\section{The MODA*-G Framework}
\label{sec:framework}
% ============================================================

\subsection{Engine I: Robust MCD Score
  ($\text{SDN}_\text{mcd}$)}
\label{sec:engine1}

\begin{definition}[MCD Estimator]
\label{def:mcd}
Given $h$ with $\lfloor n/2 \rfloor + 1 \leq h \leq n$,
the MCD estimator finds $\hat{H} \subset \{1,\ldots,n\}$
with $|\hat{H}| = h$ minimising
$\det(\hat{C}_H)$, yielding robust location
$T_n = h^{-1}\sum_{i \in \hat{H}} x_{\text{num},i}$
and scatter
$C_n = h^{-1}\sum_{i \in \hat{H}}
(x_{\text{num},i} - T_n)(x_{\text{num},i} - T_n)^\top
+ \eta I_p$, $\eta = 10^{-8}$.
\end{definition}

\begin{definition}[$\text{SDN}_\text{mcd}$ Score]
\label{def:smcd}
$D_\text{mcd}(x) = \sqrt{(x_\text{num} - T_n)^\top
C_n^{-1}(x_\text{num} - T_n)}$,\quad
$\text{SDN}_\text{mcd}(x) = \varphi(D_\text{mcd}(x);
\kappa)$, $\kappa = \sqrt{\chi^2_{p,0.975}}$.
\end{definition}

The C-steps algorithm of \citet{rousseeuw1999}
approximates the global minimum. The breakdown point is
$(n-h+1)/n \to 50\%$ \citep{rousseeuw1985}.

\subsection{Engine II: Robust Pe\~{n}a Kurtosis Score
  ($\text{SDN}_\text{pe\~{n}a}$)}
\label{sec:engine2}

\begin{definition}[Extreme kurtosis directions]
\label{def:Vstar}
Let $V^* = V^*_{\max} \cup V^*_{\min}$ denote the $2p$
directions maximising and minimising the projection
kurtosis $\hat{\kappa}(v) = \E[(v^\top x)^4] /
[\E[(v^\top x)^2]]^2$ over $\mathbb{S}^{p-1}$,
computed iteratively following \citet{pena2001}.
\end{definition}

\begin{definition}[$\text{SDN}_\text{pe\~{n}a}$ Score
  --- \emph{Original contribution}]
\label{def:spena}
$r_v(x) = |v^\top x_\text{num} - \hat{\mu}_v|/
\hat{\sigma}_v$, where
$\hat{\mu}_v = \median\{v^\top x_i\}$,
$\hat{\sigma}_v = 1.4826\cdot\MAD\{v^\top x_i\}$.
\[
  \text{SDN}_\text{pe\~{n}a}(x) =
  \varphi\!\left(\max_{v \in V^*} r_v(x);\; 3\right).
\]
\end{definition}

\begin{remark}
The original Curtosis-1 \citep{pena2001} evaluates
projection outlyingness using classical moments, giving
breakdown point $\varepsilon^* = 0$. MODA*-G replaces
this with median/MAD evaluation, raising
$\varepsilon^*(\text{SDN}_\text{pe\~{n}a})$ from 0\%
to 50\%---an original contribution established formally
in Theorem~\ref{thm:breakdown}.
\end{remark}

\subsection{Engine III: Gower-MAD Distance Score
  ($\text{SDG}_\text{mad}$)}
\label{sec:engine3}

\begin{definition}[Gower-MAD distance
  --- \emph{Original contribution}]
\label{def:gower}
For $x, y \in \X$:
\[
  d_G(x, y) = \frac{1}{p+q}\!\left[
    \sum_{j=1}^{p} \frac{|x_j - y_j|}{1.4826\cdot\MAD_j}
    + \sum_{k=1}^{q} \mathbf{1}[x_k \neq y_k]
  \right],
\]
where $\MAD_j = \MAD(\{x_{ij}\}_{i=1}^n)$.
\end{definition}

\begin{remark}
\citet{gower1971} uses $\mathrm{range}(X_j)$ as
denominator, giving $\varepsilon^*(\text{SDG}) = 0$.
Gower-MAD replaces range with $1.4826\cdot\MAD_j$,
raising $\varepsilon^*$ to 50\%
(Theorem~\ref{thm:breakdown}).
\end{remark}

\begin{definition}[$\text{SDG}_\text{mad}$ Score]
\label{def:sgower}
For $k_n = \max(3, \lfloor\sqrt{n}\rfloor)$-nearest
neighbours $\mathcal{N}_{k_n}(x)$ under $d_G$:
\[
  \bar{d}_\text{loc}(x) = k_n^{-1}
  \!\sum_{y \in \mathcal{N}_{k_n}(x)}\! d_G(x,y),
  \quad
  \bar{d}_\text{glob}(x) = (n{-}1)^{-1}
  \!\sum_{i=1}^n d_G(x, x_i),
\]
\[
  \text{SDG}_\text{mad}(x) = \varphi\!\bigl(
  5\cdot\max(\bar{d}_\text{loc}(x),\,
  0.5\cdot\bar{d}_\text{glob}(x));\;1\bigr).
\]
The constant $c = 5$ is chosen such that
$\text{SDG}_\text{mad}$ scores are comparable in
magnitude to the other engines under the null
distribution. Combining local and global distances
prevents compact outlier clusters from scoring falsely
low due to internal cohesion.
\end{definition}

\subsection{Engine IV: Categorical Entropy Score
  ($\text{SDC}$)}
\label{sec:engine4}

\begin{definition}[SDC Score --- \emph{Original
  contribution}]
\label{def:sdc}
For $c \in \C_k$ observed: $\hat{P}(X_k = c) =
n^{-1}\sum_{i=1}^n \mathbf{1}[x_{ik}=c]$.
For $c \notin \hat{\mathcal{S}}_k$ (unobserved):
$\hat{P}_\text{unseen} = 1/(n + m_k + 1)$.
\[
  \text{SDC}(x) = 1 - \left[\prod_{k=1}^{q}
  \hat{P}(X_k = x_k)\right]^{1/q}.
\]
\end{definition}

\subsection{Softmax Gating}
\label{sec:gating}

\begin{definition}[Softmax Gating weights]
\label{def:gating}
Let $\mathbf{S}(x) = (S_1(x), \ldots, S_4(x))^\top$
where $S_1 = \text{SDN}_\text{mcd}$,
$S_2 = \text{SDN}_\text{pe\~{n}a}$,
$S_3 = \text{SDG}_\text{mad}$, $S_4 = \text{SDC}$.
For temperature $k > 0$:
\begin{equation}
  w_i(x; k) = \frac{\exp(k \cdot S_i(x))}
  {\sum_{j=1}^{4} \exp(k \cdot S_j(x))},
  \quad i = 1,2,3,4.
  \label{eq:softmax}
\end{equation}
\end{definition}

\begin{definition}[MODA*-G Score and Threshold]
\label{def:modag}
\begin{equation}
  \moda(x; k) = \sum_{i=1}^{4} w_i(x; k) \cdot S_i(x),
  \label{eq:modag}
\end{equation}
\[
  \tau^* = \median_i\{\moda(x_i;k)\}
           + 3\cdot\MAD_i\{\moda(x_i;k)\}.
\]
Observation $x$ is classified anomalous if
$\moda(x;k) \geq \tau^*$.
\end{definition}

\begin{remark}[Limiting behaviour of $k$]
As $k \to 0$, weights converge to $1/4$ (uniform average,
equivalent to equal-weight linear combination). As
$k \to \infty$, all weight concentrates on the
highest-scoring engine (hard argmax). Our simulation
identifies $k=3$ as the optimal balance.
\end{remark}

\begin{remark}[Numerical stability]
Equation~\eqref{eq:softmax} is computed via the
log-sum-exp trick:
$w_i(x;k) = \exp(kS_i - \max_j kS_j) /
\sum_j \exp(kS_j - \max_j kS_j)$.
\end{remark}

\begin{algorithm}[H]
\caption{MODA*-G: Complete Algorithm}
\label{alg:modag}
\begin{algorithmic}[1]
\Require $\mathbf{X}_n \subset \R^p \times \C^q$,
  temperature $k > 0$ (default: $k=3$)
\Ensure Scores $\{\moda(x_i;k)\}$, threshold $\tau^*$,
  labels $\hat{y}_i \in \{0,1\}$
\State Compute $\text{SDN}_\text{mcd}$
  via C-steps (Def.~\ref{def:smcd})
\State Compute $V^*$; evaluate $\text{SDN}_\text{pe\~{n}a}$
  via median/MAD (Def.~\ref{def:spena})
\State Compute $d_G$ matrix; combine local and global
  distances for $\text{SDG}_\text{mad}$
  (Def.~\ref{def:sgower})
\State Estimate frequencies with Laplace smoothing;
  compute $\text{SDC}$ (Def.~\ref{def:sdc})
\State Compute Softmax weights $w_i(x_j; k)$
  (Eq.~\eqref{eq:softmax})
\State Compute $\moda(x_j; k)$ (Eq.~\eqref{eq:modag})
\State Compute $\tau^*$ (Def.~\ref{def:modag})
\State \Return scores, $\tau^*$,
  $\hat{y}_j = \mathbf{1}[\moda(x_j;k) \geq \tau^*]$
\end{algorithmic}
\end{algorithm}

\noindent\textbf{Time complexity.} The dominant cost is
$O(n^2(p+q))$ for the Gower-MAD matrix. Approximate kNN
methods reduce this to $O(n\log n)$ for large $n$.

% ============================================================
\section{Theoretical Properties}
\label{sec:theory}
% ============================================================

All results assume conditions (C1), (C2), (C3$'$), (C4)
unless stated otherwise, and hold for any $k \geq 0$.
Condition (C5) is additionally required for
Proposition~\ref{prop:rate}.

% ── Theorem 1 ─────────────────────────────────────────────
\begin{theorem}[Consistency]
\label{thm:consistency}
Under (C1), (C2), (C3$'$), for any fixed $x \in \X$:
$\moda_n(x; k) \inprob \moda(x; k)$ as $n \to \infty$.
\end{theorem}

\begin{proof}
\textbf{SDN$_\text{mcd}$:} By Theorem~1 of
\citet{rousseeuw1999}, $T_n \inprob T^*$ and
$C_n \inprob C^*$ under (C1)--(C2). Continuity of
$D_\text{mcd}^2(x; T, C)$ and $\varphi$ give
$\text{SDN}_\text{mcd,n}(x) \inprob
\text{SDN}_\text{mcd}(x)$ by the continuous mapping
theorem.

\textbf{SDN$_\text{pe\~{n}a}$:} By Glivenko-Cantelli
for the class $\{x \mapsto v^\top x : v \in
\mathbb{S}^{p-1}\}$, $\hat{\kappa}(v) \to \kappa(v)$
uniformly. Uniform convergence implies $v_{k,n} \inprob
v_k^*$. Consistency of median and MAD gives
$\text{SDN}_\text{pe\~{n}a,n}(x) \inprob
\text{SDN}_\text{pe\~{n}a}(x)$.

\textbf{SDG$_\text{mad}$:} $\MAD_{j,n} \inprob \MAD_j$
by the SLLN. kNN consistency under $k_n \to \infty$,
$k_n/n \to 0$ follows from \citet{devroye1996}. Hence
$\text{SDG}_\text{mad,n}(x) \inprob
\text{SDG}_\text{mad}(x)$.

\textbf{SDC:} Under (C3$'$), all category frequencies
satisfy $P(X_k = c) > 0$ under the mixture $P$, so
$\hat{P}_n(X_k = c) \inprob P(X_k = c)$ by the SLLN
for all $c$ including those absent from $P_0$. The
geometric mean is continuous in the probabilities.
Hence $\text{SDC}_n(x) \inprob \text{SDC}(x)$.

\textbf{Softmax:} The map $\mathbf{s} \mapsto
\exp(ks_i)/\sum_j \exp(ks_j)$ is continuous for all
finite $k$. Componentwise convergence of engine scores
implies $w_{i,n}(x;k) \inprob w_i(x;k)$ by the
continuous mapping theorem. The product
$w_{i,n} \cdot S_{i,n}$ converges in probability, and
the finite sum $\moda_n(x;k) \inprob \moda(x;k)$.
\end{proof}

% ── Theorem 2 ─────────────────────────────────────────────
\begin{theorem}[50\% Breakdown Point]
\label{thm:breakdown}
Under (C1), (C2), (C3$'$) with Gower-MAD:
$\varepsilon^*(\moda, \mathbf{X}_n) =
(n-h+1)/n \to 50\%$ as $n \to \infty$.
\end{theorem}

\begin{proof}
\textbf{$\varepsilon^*(\text{SDN}_\text{mcd}) = 50\%$:}
MCD with $h = \lfloor 0.75n \rfloor$ achieves
$(n-h+1)/n \to 50\%$ \citep{rousseeuw1985}.

\textbf{$\varepsilon^*(\text{SDN}_\text{pe\~{n}a})
= 50\%$:}
Classical Curtosis-1 uses moments with
$\varepsilon^* = 0$ \citep{pena2001}. MODA*-G replaces
with median ($\varepsilon^* = 50\%$) and MAD
($\varepsilon^* = 50\%$) \citep{rousseeuw2018}.
The composition achieves $\varepsilon^* = 50\%$.
This is an original contribution.

\textbf{$\varepsilon^*(\text{SDG}_\text{mad}) = 50\%$:}
Classical Gower uses $\mathrm{range}(X_j)$ with
$\varepsilon^* = 0$. Gower-MAD replaces range with
$1.4826\cdot\MAD_j$ ($\varepsilon^* = 50\%$). Categorical
contributions $\mathbf{1}[x_k \neq y_k] \in \{0,1\}$
are bounded. Hence $\varepsilon^* = 50\%$.
This is an original contribution.

\textbf{$\varepsilon^*(\text{SDC}) = 50\%$:}
SDC $\in [0,1)$ by construction regardless of
contamination.

\textbf{Gating preserves breakdown point:}
Softmax weights $w_i(x;k)$ are continuous functions of
bounded engine scores. Under contamination
$\varepsilon < 50\%$, the clean majority dominates the
score distribution, and no corrupted engine can
systematically bias the gated combination beyond the
breakdown threshold. Therefore:
$\varepsilon^*(\moda) = \min\{50\%, 50\%, 50\%, 50\%\}
= 50\%$.
\end{proof}

\begin{remark}
Without Gower-MAD---using the classical range---the
breakdown point of MODA*-G would be 0\%, dominated by
SDG. The Gower-MAD robustification is necessary, not
merely a refinement, to achieve $\varepsilon^* = 50\%$.
\end{remark}

% ── Theorem 3 ─────────────────────────────────────────────
\begin{theorem}[Partial Affine Equivariance]
\label{thm:equivariance}
Let $T_A(x_\text{num}, x_\text{cat}) =
(Ax_\text{num} + b, x_\text{cat})$ with $A$ invertible.
Under (C1)--(C4):
(i) $\text{SDN}_\text{mcd}(T_A(x)) =
\text{SDN}_\text{mcd}(x)$ (full);
(ii) $\text{SDN}_\text{pe\~{n}a}(T_A(x)) =
\text{SDN}_\text{pe\~{n}a}(x)$ (full);
(iii) $\text{SDG}_\text{mad}(T_A(x)) =
\text{SDG}_\text{mad}(x)$ iff $A = \lambda I$;
(iv) $\text{SDC}(T_A(x)) = \text{SDC}(x)$ (invariant);
(v) $\moda(T_{\lambda I}(x);k) = \moda(x;k)$.
\end{theorem}

\begin{proof}
(i): MCD satisfies
$T_n(T_A(\mathbf{X})) = AT_n + b$ and
$C_n(T_A(\mathbf{X})) = AC_nA^\top$
\citep{rousseeuw1985}.
For $y = Ax + b$:
$D_\text{mcd}^2(y) = (A(x-T_n))^\top
A^{-\top}C_n^{-1}A^{-1}A(x-T_n) = D_\text{mcd}^2(x)$.

(ii): Kurtosis is affinely invariant by Lemma~1 of
\citet{pena2001}: $\hat{\kappa}(v^\top(Ax+b)) =
\hat{\kappa}(w^\top x)$ where $w = A^\top v/\|A^\top v\|$.
Extreme directions transform consistently; median/MAD
are equivariant.

(iii): For $A = \lambda I$:
$|{\lambda x_j - \lambda y_j}|/(1.4826|\lambda|\MAD_j) =
|x_j - y_j|/(1.4826\MAD_j)$. For general $A$,
off-diagonal entries break this equality.

(iv): SDC depends only on $x_\text{cat}$, unaffected
by $T_A$.

(v): Under $A = \lambda I$, all $S_i(T_{\lambda I}(x)) =
S_i(x)$ by (i)--(iv), so $w_i$ and $\moda$ are
invariant.
\end{proof}

\begin{remark}
Partial equivariance of SDG$_\text{mad}$ is a structural
property of all Gower-type distances and does not affect
utility in tabular settings where variables have fixed,
interpretable units.
\end{remark}

% ── Theorem 4 ─────────────────────────────────────────────
\begin{theorem}[Score Comparability]
\label{thm:calibration}
Under (C1), (C2), (C3$'$) and $k \geq 0$:
(i) all engine scores and $\moda(x;k) \in [0,1)$;
(ii) each engine score is strictly increasing in the
degree of outlyingness;
(iii) $\E[S_i(x)] < 0.5$ for clean $x$ under $H_0$;
(iv) $\moda(x;k) \in [0,1)$ with consistent ordinal
interpretation.
\end{theorem}

\begin{proof}
(i): $\varphi(d;\kappa) \in [0,1)$ for all finite $d$.
SDC $\in [0,1)$ under (C3$'$). Convex combination
of $[0,1)$ values is in $[0,1)$.
(ii): Each score is a composition of increasing functions
in outlyingness; the Softmax preserves monotonicity by
amplifying higher scores.
(iii): Under $H_0$, $D_\text{mcd}^2(x) \sim^a \chi^2_p$,
giving $P(\text{SDN}_\text{mcd} < 0.5 \mid H_0) \approx
0.975$. Analogous for other engines.
(iv): Follows from (i) and (iii).
\end{proof}

\begin{corollary}[False positive rate control]
\label{cor:fpr}
$\tau^* = \median(\moda) + 3\cdot\MAD(\moda)$ satisfies
$P(\moda(x;k) \geq \tau^* \mid H_0) \to \alpha_\text{FP}
\approx 0.3\%$ asymptotically.
\end{corollary}

% ── Theorem 5 ─────────────────────────────────────────────
\begin{theorem}[Stochastic Dominance]
\label{thm:dominance}
Define $\beta(S, \tau, x^*) = P(S(x^*) \geq \tau \mid
x^* \text{ is an outlier})$. For each canonical type
$O_i$ ($i=1,2,3,4$), there exists at least one engine
$M_j$ such that for any $k \geq 0$:
$\beta(M_j, \tau, x^*) < \beta(\moda(\cdot;k),
\tau^*, x^*)$.
\end{theorem}

\begin{proof}
\textbf{O1 (isolated numerical):}
$\text{SDC}(x^*) \approx 0$ (typical categories), so
$\beta(\text{SDC}) \approx 0$. With
$S_\text{mcd} \gg S_\text{sdc}$, Softmax assigns
$w_\text{mcd} \to 1$ and $\moda(x^*;k) \to 1$.

\textbf{O2 (clustered numerical):}
MCD suffers masking, $\beta(\text{SDN}_\text{mcd})
\approx 0$ \citep{pena2001}. Pe\~{n}a's minimum-kurtosis
direction detects the cluster; Softmax amplifies
$w_\text{pe\~{n}a} \to 1$.

\textbf{O3 (moderate mixed):}
MCD and Pe\~{n}a individually fail; SDG$_\text{mad}$
and SDC jointly detect. Softmax amplifies whichever
scores higher for each observation.

\textbf{O4 (pure categorical):}
$\beta(\text{SDN}_\text{mcd}) = \beta(\text{SDN}
_\text{pe\~{n}a}) \approx 0$.
Under (C3$'$), $\hat{P}(c_{99}) \to
\varepsilon P_\text{out}(c_{99})/n \approx 0$ for small
$\varepsilon$, making $\text{SDC}(x^*) \to 1$.
Softmax with $k=3$ assigns $w_\text{sdc} \geq 0.84$
(Proposition~\ref{prop:gating}), giving
$\moda(x^*;k) \to 1$. Simulation confirms
AUC $= 1.000$ at $\varepsilon \leq 0.05$.
\end{proof}

\begin{corollary}[Necessity of all four engines]
\label{cor:necessity}
No proper subset of the four engines covers all four
canonical types: SDC fails on $O_1$; SDN$_\text{mcd}$
fails on $O_2$, $O_3$, $O_4$; SDN$_\text{pe\~{n}a}$
fails on $O_3$, $O_4$; SDG$_\text{mad}$ alone fails
on $O_4$. All four engines are necessary.
\end{corollary}

% ── Proposition 1 ─────────────────────────────────────────
\begin{proposition}[Softmax Gating Dominance]
\label{prop:gating}
If $S_{k^*}(x^*) \geq S_j(x^*) + \delta$ for all
$j \neq k^*$ and some $\delta > 0$, then:
\begin{equation}
  w_{k^*}(x^*; k) \geq
  \frac{e^{k\delta}}{3 + e^{k\delta}} \to 1
  \quad \text{as } k \to \infty.
  \label{eq:prop1}
\end{equation}
\end{proposition}

\begin{proof}
$w_{k^*} = e^{kS_{k^*}}/\sum_j e^{kS_j}
\geq e^{kS_{k^*}} / (e^{kS_{k^*}} +
3e^{k(S_{k^*}-\delta)}) =
1/(1 + 3e^{-k\delta}) = e^{k\delta}/(3+e^{k\delta})$.
Since $e^{k\delta} \to \infty$, $w_{k^*} \to 1$,
and $\moda(x^*;k) \geq w_{k^*} S_{k^*} \to S_{k^*}$.
\end{proof}

\begin{remark}
In the O4 scenario with $S_\text{sdc}(x^*)=0.947$,
$S_\text{mcd}(x^*)=0.023$, and $k=3$:
$w_\text{sdc} \geq e^{3 \times 0.924}/
(3 + e^{3 \times 0.924}) = 15.9/18.9 = 0.841$,
so SDC receives 84.1\% of total weight automatically.
\end{remark}

% ── Proposition 2 ─────────────────────────────────────────
\begin{proposition}[$O_P(n^{-1/2})$ Convergence Rate]
\label{prop:rate}
Under (C1), (C2), (C3$'$), (C4), (C5) and fixed
$k \geq 0$:
\begin{enumerate}[label=(\roman*)]
  \item \emph{Rate for each engine:}
    $S_{i,n}(x^*) - S_i(x^*) = O_P(n^{-1/2})$
    for $i = 1,2,3,4$.
  \item \emph{Rate for gating weights:}
    $w_{i,n}(x^*;k) - w_i(x^*;k) = O_P(n^{-1/2})$.
  \item \emph{Rate for MODA*-G score:}
    $|\moda_n(x^*;k) - \moda(x^*;k)| =
    O_P(n^{-1/2})$.
  \item \emph{Asymptotic normality:}
    $\sqrt{n}(\moda_n(x^*;k) - \moda(x^*;k))
    \indist \mathcal{N}(0, \sigma^2_\moda(x^*;k))$,
    where $\sigma^2_\moda = (\nabla h)^\top \Sigma
    (\nabla h)$, $h(\mathbf{S}) = \sum_i g_i(\mathbf{S})
    S_i$, and $\Sigma$ is the joint asymptotic covariance
    of $(S_{1,n}, S_{2,n}, S_{3,n}, S_{4,n})$.
\end{enumerate}
\end{proposition}

\begin{proof}
\textbf{Part (i):}
For SDN$_\text{mcd}$: the MCD estimators satisfy
$\sqrt{n}(T_n - T^*) \indist \mathcal{N}(0, \Sigma_T)$
and $\sqrt{n}\,\mathrm{vec}(C_n - C^*) \indist
\mathcal{N}(0, \Sigma_C)$ under (C1), (C2), (C5)
\citep{rousseeuw1999}. The delta method applied to the
continuous differentiable map
$(T, C) \mapsto \varphi(D_\text{mcd}(x; T, C))$ gives
$S_\text{mcd,n}(x^*) - S_\text{mcd}(x^*) =
O_P(n^{-1/2})$.
For SDN$_\text{pe\~{n}a}$: empirical kurtosis
$\hat{\kappa}(v)$ converges at rate $n^{-1/2}$
uniformly over $\mathbb{S}^{p-1}$ under (C5). The
median and MAD are $n^{-1/2}$-consistent. The delta
method gives the same rate.
For SDG$_\text{mad}$: $\MAD_{j,n} - \MAD_j =
O_P(n^{-1/2})$ under (C5) \citep{rousseeuw2018}.
kNN consistency with $k_n = O(\sqrt{n})$ gives
$O_P(n^{-1/2})$ for the kNN mean distance.
For SDC: by the CLT applied to the i.i.d.\ indicators
$\mathbf{1}[x_{ik} = c]$, each frequency satisfies
$\hat{P}_n(X_k = c) - P(X_k = c) = O_P(n^{-1/2})$.
The geometric mean is differentiable in frequencies
(under $P(X_k = c) > 0$ guaranteed by (C3$'$) for all
$c$ under the mixture), so the delta method applies.

\textbf{Part (ii):}
The Softmax $g_i : \R^4 \to \R$,
$g_i(\mathbf{s}) = e^{ks_i}/\sum_j e^{ks_j}$,
is infinitely differentiable with bounded partial
derivatives:
$\partial g_i/\partial s_i = k w_i(1 - w_i)$ and
$\partial g_i/\partial s_j = -k w_i w_j$ ($j \neq i$).
By the multivariate delta method applied to the
composition $g_i \circ \mathbf{S}_n$:
$\sqrt{n}(g_i(\mathbf{S}_n) - g_i(\mathbf{S}))
\indist \mathcal{N}(0, (\nabla g_i)^\top \Sigma
\nabla g_i)$,
where $\Sigma$ is the joint asymptotic covariance of
the four engine scores (finite under (C5)).
Hence $w_{i,n}(x^*;k) - w_i(x^*;k) = O_P(n^{-1/2})$.

\textbf{Part (iii):}
Write $\moda_n - \moda = A_n + B_n$ where:
$A_n = \sum_i (w_{i,n} - w_i) S_{i,n}
= O_P(n^{-1/2})$
since $w_{i,n} - w_i = O_P(n^{-1/2})$ and $S_{i,n}
\in [0,1)$ is bounded;
$B_n = \sum_i w_i (S_{i,n} - S_i)
= O_P(n^{-1/2})$
since $w_i \in [0,1]$ and $S_{i,n} - S_i =
O_P(n^{-1/2})$ from (i).
Therefore $|\moda_n - \moda| \leq |A_n| + |B_n| =
O_P(n^{-1/2})$.

\textbf{Part (iv):}
Define $h : \R^4 \to \R$ by $h(\mathbf{s}) =
\sum_i g_i(\mathbf{s}) s_i$, which is continuously
differentiable. By the multivariate delta method:
$\sqrt{n}(h(\mathbf{S}_n) - h(\mathbf{S}))
\indist \mathcal{N}(0, (\nabla h)^\top \Sigma
\nabla h)$,
establishing asymptotic normality with variance
$\sigma^2_\moda = (\nabla h)^\top \Sigma \nabla h
< \infty$ under (C5).
\end{proof}

\begin{remark}
The $O_P(n^{-1/2})$ rate implies that with $n = 300$,
the typical estimation error in the MODA*-G score is
of order $\sigma_\moda/\sqrt{300} \approx \sigma_\moda/
17.3$. In our simulation, the score separation between
outliers and normal observations in O4 is approximately
0.924; the estimation error is three orders of magnitude
smaller, confirming the reliability of $n = 300$ for
the simulation study.
\end{remark}

% ============================================================
\section{Simulation Study}
\label{sec:simulation}
% ============================================================

\subsection{Design}

Monte Carlo simulation: $n = 300$, $p = 4$, $q = 2$
(3 categories each), 100 replications per cell,
$\varepsilon \in \{0.01, 0.05, 0.10, 0.15, 0.20\}$.
Clean: $X_\text{num} \sim \mathcal{N}_4(\mathbf{0},
I_4)$, $X_\text{cat} \sim \text{Multinomial}
(0.6, 0.3, 0.1)^2$.
Tested: MODA*-G with $k \in \{1,2,3,5,8,10\}$,
linear base ($k \to 0$), Isolation Forest
\citep{liu2008}, LOF \citep{breunig2000} (both with
label-encoded categoricals). Metric: AUC-ROC
\citep{fawcett2006}.

\subsection{Canonical Outlier Types}

\begin{description}[leftmargin=*]
\item[$O_1$: Isolated numerical] $x^* = (\mu + 6v,
  c_\text{typ})$, $v \in \mathbb{S}^{p-1}$ random,
  $\text{Var} = 0.1$. Typical categories.
\item[$O_2$: Clustered numerical] $m$ compact outliers
  at $4\sigma$ from origin, $\text{Var} = 0.05$,
  typical categories. Causes MCD masking.
\item[$O_3$: Moderate mixed] $x^* = (\mu + 2v, c_3)$
  where $c_3$ is a rarely-occurring category.
\item[$O_4$: Pure categorical] $x^* = (\mu', c_{99})$,
  $\mu' \sim \mathcal{N}_4(\mathbf{0}, 0.1I_4)$,
  $c_{99}$ never observed in clean data.
\end{description}

\subsection{Results}

\begin{table}[ht]
\centering
\caption{Mean AUC-ROC ($\varepsilon \leq 0.15$,
  100 reps). Best result per scenario in \textbf{bold}.
  MODA*-G ($k{=}3$) outperforms all competitors in
  all four scenarios simultaneously.}
\label{tab:main}
\small
\begin{tabular}{lcccccc}
\toprule
Scenario & MODA*-G & MODA* & IF & LOF
  & $\Delta$base & $\Delta$IF \\
 & ($k=3$) & (base) & & & & \\
\midrule
$O_1$ Isolated num.
  & $\mathbf{0.957}$ & $0.950$ & $0.903$ & $0.659$
  & $+0.008$ & $+0.054$ \\
$O_2$ Clustered num.
  & $\mathbf{0.993}$ & $0.984$ & $0.957$ & $0.670$
  & $+0.009$ & $+0.036$ \\
$O_3$ Moderate mixed
  & $\mathbf{0.986}$ & $0.936$ & $0.930$ & $0.623$
  & $+0.050$ & $+0.056$ \\
$O_4$ Pure categorical
  & $\mathbf{0.977}$ & $0.513$ & $0.855$ & $0.577$
  & $+0.464$ & $+0.122$ \\
\midrule
\textbf{Mean}
  & $\mathbf{0.978}$ & $0.846$ & $0.911$ & $0.623$
  & $+0.132$ & $+0.067$ \\
\bottomrule
\end{tabular}
\end{table}

\begin{table}[ht]
\centering
\caption{O4: AUC-ROC by $\varepsilon$ and temperature
  $k$. MODA*-G with $k \geq 3$ achieves AUC $= 1.000$
  at $\varepsilon \leq 0.05$. Linear base fails.}
\label{tab:o4}
\small
\begin{tabular}{lcccccccc}
\toprule
$\varepsilon$ & Base & $k{=}1$ & $k{=}2$ & $k{=}3$
  & $k{=}5$ & $k{=}8$ & $k{=}10$ & SDC only\\
\midrule
$0.01$ & $0.706$ & $0.947$ & $1.000$ & $\mathbf{1.000}$
  & $1.000$ & $1.000$ & $1.000$ & $1.000$ \\
$0.05$ & $0.596$ & $0.859$ & $0.992$ & $\mathbf{1.000}$
  & $1.000$ & $1.000$ & $1.000$ & $1.000$ \\
$0.10$ & $0.446$ & $0.702$ & $0.943$ & $0.988$
  & $\mathbf{0.996}$ & $0.995$ & $0.994$ & $0.992$ \\
$0.15$ & $0.302$ & $0.495$ & $0.802$ & $0.919$
  & $0.966$ & $\mathbf{0.967}$ & $0.965$ & $0.956$ \\
$0.20$ & $0.172$ & $0.288$ & $0.572$ & $0.747$
  & $0.856$ & $\mathbf{0.859}$ & $0.848$ & $0.867$ \\
\bottomrule
\end{tabular}
\end{table}

\begin{table}[ht]
\centering
\caption{Temperature sensitivity: mean AUC-ROC
  ($\varepsilon \leq 0.15$). $k=3$ achieves the
  highest mean (0.978).}
\label{tab:sensitivity}
\small
\begin{tabular}{lccccc}
\toprule
$k$ & $O_1$ & $O_2$ & $O_3$ & $O_4$ & Mean \\
\midrule
base & $0.950$ & $0.984$ & $0.936$ & $0.513$ & $0.846$\\
$1$  & $0.980$ & $0.997$ & $0.959$ & $0.751$ & $0.921$\\
$2$  & $0.970$ & $0.996$ & $0.980$ & $0.934$ & $0.970$\\
$\mathbf{3}$
     & $\mathbf{0.957}$ & $\mathbf{0.993}$
     & $\mathbf{0.986}$ & $\mathbf{0.977}$
     & $\mathbf{0.978}$\\
$5$  & $0.933$ & $0.983$ & $0.987$ & $0.990$ & $0.973$\\
$8$  & $0.924$ & $0.977$ & $0.987$ & $0.991$ & $0.970$\\
$10$ & $0.928$ & $0.978$ & $0.987$ & $0.990$ & $0.971$\\
\bottomrule
\end{tabular}
\end{table}

\begin{table}[ht]
\centering
\caption{Individual engine AUC-ROC
  ($\varepsilon \leq 0.15$). No single engine achieves
  strong performance across all scenarios, confirming
  Corollary~\ref{cor:necessity}.}
\label{tab:engines}
\small
\begin{tabular}{lcccc}
\toprule
Engine & $O_1$ & $O_2$ & $O_3$ & $O_4$ \\
\midrule
$\text{SDN}_\text{mcd}$   & $0.960$
  & $0.994$ & $0.593$ & $0.026$ \\
$\text{SDN}_\text{pe\~{n}a}$ & $0.999$
  & $1.000$ & $0.614$ & $0.026$ \\
$\text{SDG}_\text{mad}$   & $0.996$
  & $1.000$ & $0.825$ & $0.345$ \\
$\text{SDC}$               & $0.497$
  & $0.500$ & $0.987$ & $0.987$ \\
\midrule
MODA*-G ($k=3$) & $\mathbf{0.957}$
  & $\mathbf{0.993}$ & $\mathbf{0.986}$
  & $\mathbf{0.977}$ \\
\bottomrule
\end{tabular}
\end{table}

\paragraph{Discussion.}
MODA*-G with $k=3$ achieves AUC $> 0.95$ across all
four canonical scenarios at $\varepsilon \leq 0.15$,
outperforming IF in all scenarios simultaneously for
the first time. The O4 breakthrough ($+46.4\%$ vs.\
linear base, $+12.2\%$ vs.\ IF) validates
Proposition~\ref{prop:gating}: when SDC scores 0.947
and MCD scores 0.023, the Softmax assigns
$w_\text{sdc} \approx 0.841$ automatically. The
temperature sensitivity analysis confirms robustness:
mean AUC $> 0.97$ for all $k \geq 2$, ruling out
parameter-specific overfitting. For applications where
pure categorical contamination predominates, $k=5$
maximises O4 (AUC $= 0.990$) at modest cost to O1
(AUC $= 0.933$).

% ============================================================
\section{Real Data Study --- Annthyroid Benchmark}
\label{sec:realdata}
% ============================================================

\subsection{Dataset and Protocol}

We evaluate MODA*-G on the Annthyroid dataset
\citep{campos2016}, the standard benchmark for
mixed-data outlier detection established by
\citet{campos2016}. The dataset contains clinical
measurements for patients evaluated for thyroid
disorders, with three classes: normal (inlier),
hyperfunction (outlier), and subnormal functioning
(outlier).

\textbf{Data version:} We use the pre-normalised
version from the LMU Munich benchmark repository
\citep{campos2016}, file \texttt{Annthyroid\_07.arff},
which scales all numerical features to $[0,1]$ using
range normalisation. This is the standard version
used in all benchmark comparisons in the outlier
detection literature; using raw UCI data without
normalisation degrades all distance-based methods
due to the extreme scale difference between features
(e.g.\ TSH ranges $[0,530]$ mU/L vs.\ T4U $[0.2,2.3]$).

\begin{table}[ht]
\centering
\caption{Annthyroid dataset characteristics
  (Campos et al.\ 2016 benchmark version).}
\label{tab:dataset}
\small
\begin{tabular}{ll}
\toprule
Property & Value \\
\midrule
Total observations ($n$) & 7,200 \\
Normal (inlier) & 6,666 (92.58\%) \\
Outliers & 534 (7.42\%) \\
Numerical features ($p$) & 6 (age, TSH, T3, TT4, T4U, FTI) \\
Categorical features ($q$) & 15 (binary clinical symptoms) \\
Total features & 21 \\
Outlier definition & Hyperfunction + subnormal classes \\
\bottomrule
\end{tabular}
\end{table}

\textbf{Evaluation protocol:} Following
\citet{campos2016}, we report AUC-ROC as the primary
metric. We use 50 stratified bootstrap replicates
with subsampling of $n_\text{sub} = 2{,}000$
observations per replicate, preserving the 7.42\%
outlier proportion. Standard deviations are reported
across bootstrap replicates.

\subsection{Competing Methods}

We compare MODA*-G against six methods spanning the
major paradigms for unsupervised outlier detection:

\begin{description}[leftmargin=*]
\item[IF] Isolation Forest \citep{liu2008}: tree-based
  path-length isolation. Label-encoded categorical features.
\item[EIF] Extended Isolation Forest
  \citep{hariri2019}: random hyperplane cuts replacing
  axis-parallel splits. Same encoding as IF.
\item[COPOD] Copula-Based Outlier Detection
  \citep{li2020copod}: empirical copula tail probabilities
  per dimension, aggregated via log-sum. Requires
  continuous features; categorical features are
  label-encoded.
\item[CBOD] CatBoost-style Outlier Detection
  (unsupervised): gradient-boosted reconstruction
  residuals with target-encoded categorical features,
  following the knowledge-uncertainty principle of
  \citet{malinin2020}.
\item[LOF] Local Outlier Factor \citep{breunig2000}:
  density-based local outlyingness.
  Label-encoded categorical features.
\item[MODA*-G] Our method: native operation in
  $\R^p \times \C^q$, $k=3$. No encoding performed.
\end{description}

\subsection{Results}

\begin{table}[ht]
\centering
\caption{Real-data benchmark --- Annthyroid\_07.arff.
  AUC-ROC (mean $\pm$ std, 50 bootstrap replicates).
  MODA*-G outperforms all six competitors.
  Methods marked $\dagger$ use label encoding of
  categorical features.}
\label{tab:realdata}
\small
\begin{tabular}{lccc}
\toprule
Method & AUC-ROC & $\pm$ & $\Delta$ vs IF \\
\midrule
\textbf{MODA*-G ($k=3$)} & $\mathbf{0.939}$ & $0.010$
  & $+0.307$ \\
IF$^\dagger$ & $0.632$ & $0.021$ & --- \\
LOF$^\dagger$ & $0.620$ & $0.023$ & $-0.012$ \\
CBOD$^\dagger$ & $0.554$ & $0.024$ & $-0.078$ \\
COPOD$^\dagger$ & $0.553$ & $0.094$ & $-0.079$ \\
EIF$^\dagger$ & $0.496$ & $0.064$ & $-0.136$ \\
\midrule
Best method, \citet{campos2016} & ${\sim}0.75$ &
  --- & --- \\
\bottomrule
\end{tabular}
\end{table}

\subsection{Discussion}

\paragraph{MODA*-G dominance.}
MODA*-G achieves AUC $= 0.939$, outperforming the
best competitor (IF, $0.632$) by $+0.307$ AUC points.
This advantage is attributable to native mixed-data
operation: the 15 binary categorical symptom features
are processed by SDC (categorical entropy score) and
Gower-MAD without any encoding, preserving the
categorical metric structure. All five competitors
encode these features as integers, distorting the
metric.

\paragraph{EIF underperforms IF.}
The EIF result (AUC $= 0.496$, near-random) illustrates
an important property of EIF: random hyperplane cuts
improve over axis-parallel cuts when features have
heterogeneous scales. On the pre-normalised Annthyroid
data (all features in $[0,1]$), axis-parallel cuts are
already geometrically appropriate, and randomisation
introduces unnecessary variance.

\paragraph{COPOD and CBOD near-random.}
COPOD (AUC $= 0.553$) is designed for continuous
marginal distributions. The 15 binary features produce
degenerate empirical CDFs (step functions at 0 and 1),
for which the copula tail probability is uninformative.
CBOD (AUC $= 0.554$) similarly suffers from the
encoding distortion of binary-to-integer conversion.

\paragraph{Comparison with Campos et al.\ (2016).}
\citet{campos2016} benchmark density-based methods
(LOF, kNN, INFLO, COF, ODIN) on the same dataset.
Their best result is INFLO with AUC $\approx 0.75$.
MODA*-G achieves AUC $= 0.939$, surpassing INFLO
by approximately $+0.189$ AUC points. All methods in
\citet{campos2016} operate exclusively on numerical
features after encoding; none provides a dedicated
categorical anomaly engine equivalent to SDC. This
confirms that the performance gap is architectural ---
native mixed-data operation provides a systematic
advantage that cannot be recovered by encoding
categorical features as numerical values.

% ============================================================
\section{Biosurveillance Applications}
\label{sec:biosurveillance}
% ============================================================

The four canonical scenarios map directly to operational
biosurveillance challenges. O1 (isolated numerical)
corresponds to sudden spikes in case counts or laboratory
values; O2 (clustered numerical) to geographic outbreak
clusters that cause masking in standard methods; O3
(moderate mixed) to emerging pathogen signals with
moderate numerical elevation and unusual ICD-10 codes;
O4 (pure categorical) to novel pathogen detection where
a patient carries a pathogen code never previously
recorded. In this last case---the most critical for
biosecurity---MODA*-G achieves AUC $= 1.000$ at low
contamination rates, a capability no existing method
provides with formal guarantees.

The per-engine diagnostic profile routes each alert
to the appropriate response: MCD-dominant alerts suggest
measurement anomalies; Pe\~{n}a-dominant alerts suggest
outbreak clusters; SDC-dominant alerts trigger immediate
novel pathogen protocols. This aligns with priorities of
the National Biosurveillance Strategy \citep{hhs2022}
and CDC BioSense Platform requirements for interpretable,
real-time mixed-data anomaly detection.

% ============================================================
\section{Limitations}
\label{sec:limitations}
% ============================================================

We document three limitations that bound the scope of the
current results and suggest directions for future work.

\paragraph{Computational complexity.}
The dominant computational cost of MODA*-G is the
pairwise Gower-MAD distance matrix, computed in
$O(n^2(p+q))$ time and $O(n^2)$ space. For datasets
with $n > 10{,}000$, this becomes prohibitive without
approximation. Approximate kNN methods---including ball
trees, KD-trees adapted for mixed data, and random
projection approximations---can reduce this to
$O(n \log n)$ without substantially affecting detection
accuracy. Implementation of these approximations is
deferred to future work.

\paragraph{Theoretical scope under (C3$'$).}
Condition~(C3$'$) extends the standard observability
assumption to allow categories absent from the clean
distribution $P_0$ by working under the mixture model
$P = (1-\varepsilon)P_0 + \varepsilon P_\text{out}$.
The formal consistency results (Theorem~\ref{thm:consistency}
and Proposition~\ref{prop:rate}) apply when
$\varepsilon > 0$ so that all categories have positive
probability under $P$.
At the boundary $\varepsilon \to 0$, the unseen-category
branch of SDC operates via the Laplace-smoothed estimator
$\hat{P}_\text{unseen} = 1/(n + m_k + 1) \to 0$,
which is not covered by the SLLN convergence proof.
A complete treatment for the $\varepsilon \to 0$ boundary
case---potentially using Good-Turing smoothing
\citep{good1953} or support-based formulations---is a
direction for future theoretical work. Empirically,
SDC achieves AUC $= 1.000$ for O4 at $\varepsilon
\leq 0.05$, confirming practical reliability despite
this boundary.

\paragraph{Asymptotic rate without explicit variance.}
Proposition~\ref{prop:rate} establishes the
$O_P(n^{-1/2})$ rate and asymptotic normality of
$\moda_n(x^*;k)$, but does not provide a closed-form
expression for the asymptotic variance
$\sigma^2_\moda(x^*;k)$. This quantity involves the
joint asymptotic covariance matrix $\Sigma$ of the
four engine scores, which depends on the unknown
data-generating distribution $P$ in a complex way.
Estimation of $\sigma^2_\moda$ and construction of
confidence intervals for the MODA*-G score are deferred
to future work. Practitioners requiring uncertainty
quantification can use bootstrap resampling as an
interim approach.

% ============================================================
\section{Conclusion}
\label{sec:conclusion}
% ============================================================

We have presented MODA*-G, a unified anomaly detection
framework for mixed numerical and categorical data.
Five formal theorems establish consistency, 50\%
breakdown point, partial affine equivariance, score
comparability, and stochastic dominance.
Two convergence propositions establish that the Softmax
Gating weight of the dominant engine converges to one
and that the MODA*-G score converges at rate
$O_P(n^{-1/2})$.
Monte Carlo simulation confirms that MODA*-G with $k=3$
achieves mean AUC-ROC of 0.978, outperforming Isolation
Forest in all four canonical contamination scenarios
simultaneously. On the Annthyroid real-data benchmark
\citep{campos2016}, MODA*-G achieves AUC $= 0.939$,
surpassing all six competing methods including EIF,
COPOD, CatBoost-AD, and LOF, and exceeding the best
result of \citet{campos2016} by $+0.189$ AUC points.

The central insight is that the failure of linear
combinations---demonstrated across five iterative
implementation versions---is structural: fixed weights
cannot adapt when an engine actively inverts rankings.
Softmax Gating resolves this with three lines of code
and no heuristic rules, yielding an attention mechanism
that automatically amplifies the correct detector at
the observation level. The real-data benchmark confirms
that the advantage of native mixed-data operation is
not merely theoretical: encoding categorical features
as integers systematically degrades all competing
methods on datasets with substantial categorical content.

\section*{Acknowledgements}
This work was conducted independently. The author thanks
the open-source communities behind NumPy, SciPy, and
scikit-learn, which supported the simulation study.

% ── References ────────────────────────────────────────────
\bibliographystyle{apalike}
\begin{thebibliography}{99}

\bibitem[Campos et al.(2016)]{campos2016}
Campos, G.~O., Zimek, A., Sander, J., Campello,
R.~J.~G.~B., Micenkov\'{a}, B., Schubert, E.,
Assent, I., and Houle, M.~E. (2016).
On the evaluation of unsupervised outlier detection:
measures, datasets, and an empirical study.
\textit{Data Mining and Knowledge Discovery},
30(4):891--927.

\bibitem[Hariri et al.(2019)]{hariri2019}
Hariri, S., Kind, M.~C., and Brunner, R.~J. (2019).
Extended isolation forest.
\textit{IEEE Transactions on Knowledge and Data
Engineering}, 33(4):1479--1489.

\bibitem[Li et al.(2020)]{li2020copod}
Li, Z., Zhao, Y., Botta, N., Ionescu, C., and
Hu, X. (2020). COPOD: copula-based outlier detection.
\textit{IEEE International Conference on Data Mining
(ICDM)}, pp.~1118--1123.

\bibitem[Malinin et al.(2020)]{malinin2020}
Malinin, A., Chien, J., Gales, M., Knill, K.,
Kumar, S., Rayner, P., Shafran, I., and Williams, B.
(2020). Uncertainty estimation in autoregressive
structured prediction.
\textit{arXiv preprint arXiv:2002.07650}.

\bibitem[Bahdanau et al.(2015)]{bahdanau2015}
Bahdanau, D., Cho, K., and Bengio, Y. (2015).
Neural machine translation by jointly learning to align
and translate.
\textit{International Conference on Learning
Representations (ICLR 2015)}.

\bibitem[Breunig et al.(2000)]{breunig2000}
Breunig, M.~M., Kriegel, H.-P., Ng, R.~T., and
Sander, J. (2000). LOF: identifying density-based
local outliers. \textit{ACM SIGMOD Record},
29(2):93--104.

\bibitem[Cabral et al.(2025)]{cabral2025}
Cabral, E.~F., Vinces, B.~V.~S., Silva, G.~D.~F.,
Sander, J., and Cordeiro, R.~L.~F. (2025).
Efficient outlier detection in numerical and categorical
data. \textit{Data Mining and Knowledge Discovery},
39(3):1--46.

\bibitem[Devroye et al.(1996)]{devroye1996}
Devroye, L., Gy\"{o}rfi, L., and Lugosi, G. (1996).
\textit{A Probabilistic Theory of Pattern Recognition}.
Springer, New York.

\bibitem[Fawcett(2006)]{fawcett2006}
Fawcett, T. (2006). An introduction to ROC analysis.
\textit{Pattern Recognition Letters}, 27(8):861--874.

\bibitem[Good(1953)]{good1953}
Good, I.~J. (1953). The population frequencies of
species and the estimation of population parameters.
\textit{Biometrika}, 40(3-4):237--264.

\bibitem[Gower(1971)]{gower1971}
Gower, J.~C. (1971). A general coefficient of
similarity and some of its properties.
\textit{Biometrics}, 27(4):857--871.

\bibitem[HHS(2022)]{hhs2022}
U.S.\ Department of Health and Human Services (2022).
\textit{National Biosurveillance Strategy 2022--2026}.
Washington, D.C.

\bibitem[Jordan and Jacobs(1994)]{jordan1994}
Jordan, M.~I. and Jacobs, R.~A. (1994).
Hierarchical mixtures of experts and the EM algorithm.
\textit{Neural Computation}, 6(2):181--214.

\bibitem[Liu et al.(2008)]{liu2008}
Liu, F.~T., Ting, K.~M., and Zhou, Z.-H. (2008).
Isolation forest. \textit{Proceedings of the 8th IEEE
International Conference on Data Mining}, pp.~413--422.

\bibitem[Marcano and Ferm\'{i}n(2013)]{marcano2013}
Marcano, L. and Ferm\'{i}n, W. (2013).
Comparaci\'{o}n de m\'{e}todos de detecci\'{o}n de datos
an\'{o}malos multivariantes mediante un estudio de
simulaci\'{o}n.
\textit{Saber}, 25(2):192--201.
Universidad de Oriente, Venezuela.

\bibitem[Pe\~{n}a and Prieto(2001)]{pena2001}
Pe\~{n}a, D. and Prieto, F.~J. (2001).
Multivariate outlier detection and robust covariance
matrix estimation.
\textit{Technometrics}, 43(3):286--310.

\bibitem[Rousseeuw(1985)]{rousseeuw1985}
Rousseeuw, P.~J. (1985). Multivariate estimation with
high breakdown point.
\textit{Mathematical Statistics and Applications},
8:283--297.

\bibitem[Rousseeuw and Driessen(1999)]{rousseeuw1999}
Rousseeuw, P.~J. and Van Driessen, K. (1999).
A fast algorithm for the minimum covariance determinant
estimator. \textit{Technometrics}, 41(3):212--223.

\bibitem[Rousseeuw and Hubert(2018)]{rousseeuw2018}
Rousseeuw, P.~J. and Hubert, M. (2018).
Anomaly detection by robust statistics.
\textit{WIREs Data Mining and Knowledge Discovery},
8(2):e1236.

\bibitem[Rousseeuw and Van Zomeren(1990)]{rousseeuw1990}
Rousseeuw, P.~J. and Van Zomeren, B.~C. (1990).
Unmasking multivariate outliers and leverage points.
\textit{Journal of the American Statistical
Association}, 85(411):633--651.

\bibitem[Vaswani et al.(2017)]{vaswani2017}
Vaswani, A., Shazeer, N., Parmar, N., Uszkoreit, J.,
Jones, L., Gomez, A.~N., Kaiser, L., and
Polosukhin, I. (2017). Attention is all you need.
\textit{Advances in Neural Information Processing
Systems (NeurIPS)}, 30.

\bibitem[Wang and Li(2021)]{wang2021}
Wang, Y. and Li, Y. (2021). Outlier detection based
on weighted neighbourhood information network for
mixed-valued datasets.
\textit{Information Sciences}, 564:396--415.

\bibitem[Yuan et al.(2023)]{yuan2023}
Yuan, Z. et al. (2023). Attribute-weighted outlier
detection for mixed data based on parallel mutual
information. \textit{Expert Systems with Applications}.

\end{thebibliography}

% ============================================================
\appendix
\section{Development History}
\label{app:versions}
% ============================================================

MODA*-G emerged from a systematic iterative development
process. The key insight---that fixed linear combinations
are structurally incapable of handling active ranking
inversion---required six implementation versions to
establish empirically.

\begin{description}[leftmargin=*]
\item[\textbf{v0.1--v0.2}] Four-engine framework with
  fixed weights. SDC implementation bug (hardcoded
  category 4 colliding with clean data) caused
  O4 AUC $= 0.341$. Fix: guaranteed-absent code (99)
  with proper Laplace smoothing restored
  AUC $= 1.000$ for SDC alone.
\item[\textbf{v0.3}] Adaptive base weights
  $(0.30, 0.20, 0.15, 0.35)$ and Gower local+global
  distance combination. Results: O1: 0.951, O2: 0.983,
  O3: 0.936, O4: 0.513. O4 fails because Gower
  ($\gamma{=}0.15$) and MCD ($\alpha{=}0.30$) actively
  invert rankings for compact categorical clusters.
\item[\textbf{v0.4--v0.4c}] Three variants of coherence
  weighting using Spearman correlations and separation
  criteria. Each variant improved some scenarios while
  degrading others. Root cause: all approaches relied
  on the uniform-weighted score to identify candidate
  outliers, which fails when the uniform score is
  itself dominated by inverting engines.
\item[\textbf{v0.5}] Rank-based KS informativeness
  weighting. Failed across all scenarios: converting
  scores to ranks compresses information at the
  extremes. With $\varepsilon = 0.01$ (3 outliers in
  300), the rank difference between the outlier
  (rank 300) and the nearest normal (rank 299) is
  $1/300 = 0.003$, while the raw score difference was
  0.85. Ranks discard the most relevant signal.
\item[\textbf{v0.6 = MODA*-G}] Softmax Gating.
  Three lines of code, no heuristics, no thresholds.
  Mean AUC $= 0.978$ across all four scenarios.
\end{description}

\end{document}
